\section{Preliminaries and probabilistic tools}\label{sec:preliminaries}

All graphs are finite, simple, and undirected. A \emph{stable set} is a set
of pairwise nonadjacent vertices. For a vertex set $X$, write $P[X]$ for
the induced subgraph and $P-X=P[V(P)\setminus X]$. A clique model consists
of disjoint nonempty connected vertex sets, called \emph{bags}, with an
edge between every two bags. Models lift through edge contractions and
vertex deletions, so $h$ is nonincreasing under taking minors. Whenever
a subgraph is specified by its vertices, it is induced. Adding edges to
a connected or highly connected subgraph preserves the corresponding
property, so the subgraphs produced below may be replaced by induced
subgraphs on their vertices.

A graph is $k$-connected if it has more than $k$ vertices and remains
connected after deletion of fewer than $k$ vertices. Such a graph has
minimum degree at least $k$. A separator split is a partition
$V(P)=A\mathbin{\dot\cup}Z\mathbin{\dot\cup}A'$, with $A,A'$ nonempty
and anticomplete. If $|V(P)|>k$ and $P$ is not $k$-connected, there is such
a split with $|Z|<k$. Deleting $z<k$ vertices from a $k$-connected graph
leaves a $(k-z)$-connected graph. All logarithms are natural. The null
graph has $h=\ell=\chi=f=0$. Integer parameters include minor orders,
connectivity requirements, and numbers of paths; powers bounding graph
orders need not be integers.

\subsection{Fractional coloring and established structural results}

The fractional chromatic number $f(P)$ is the minimum total weight of
nonnegative weights on the stable sets of $P$ that cover every vertex
with weight at least one. Finite linear-programming duality
\cite{galekuhntucker1951} gives
\begin{equation}\label{eq:fractional-dual}
 f(P)=\max\left\{\sum_{v\in V(P)}w(v):
 w(v)\ge0,\quad \sum_{v\in I}w(v)\le1
 \text{ for every stable set }I\right\}.
\end{equation}
The primal is feasible, using singleton stable sets, and has a finite
optimum; thus this use of duality has no feasibility or boundedness
exception. For $r>0$, $f(P)\le r$ is equivalent to a distribution on stable
sets with every vertex included with probability at least $1/r$.
Indeed, divide a fractional coloring by $r$ and put any missing mass on
the empty set; conversely multiply such a distribution by $r$.
Restriction to an induced subgraph cannot increase $f$, and adding one
vertex increases it by at most one.

\begin{theorem}[Fractional minor bound]\label{thm:fractional}
Every graph $P$ satisfies $f(P)\le4h(P)$.
\end{theorem}
This follows from the stronger bound $f(P)\le2h(P)$ of Reed and Seymour
\cite[Statement~(1.3)]{ReedSeymour1998}. We keep the factor four throughout.

\begin{theorem}[Ordinary-chromatic connectivity extraction]
\label{thm:ordinary-extraction}
If $k\ge1$ is an integer and $\chi(P)\ge7k$, then $P$ has a
$k$-connected subgraph $J$ with $\chi(J)\ge\chi(P)-6k$.
\end{theorem}
We use the form stated in \cite[Theorem~5.1]{DelcourtPostle2021},
attributed there to Gir\~ao and Narayanan. There is no minor-exclusion
hypothesis. In particular the theorem applies to auxiliary clique
blow-ups, which need not inherit any minor exclusion.

\begin{theorem}[Grouping terminals]\label{thm:grouping}
There is an absolute integer $\gamma\ge1$ such that the following holds.
In a $\gamma m$-connected graph, every partition of $m\ge1$ distinct
terminals into nonempty groups can be realized by pairwise disjoint
connected vertex sets, one containing each group.
\end{theorem}
This is a consequence of the highly linked graph theorem of
Bollob\'as and Thomason \cite{BollobasThomason1996}; we use the
case with second knitting parameter one in
\cite[Definition~5.7 and Theorem~5.8]{DelcourtPostle2021}. The constant
$\gamma$ is uniform in the number and sizes of the groups. For a group
of size two, a path within its connected set joins its terminals. We
will apply this observation to any prescribed pairing of distinct
terminals.

\subsection{Bounded sampling and concentration}

The following real-parameter form of bounded stable-set sampling is
convenient. Its integer-parameter form appears in
\cite[Lemma~4.1]{LiuLuo2026}; we include the short proof.

\begin{lemma}\label{lem:bounded}
Suppose $|V(P)|\le n$, where $n\ge1$, and $f(P)\le r$, where $r\ge1$.
There is a distribution on stable sets $I$ such that
\[
 |I|\le\lceil n/r\rceil,
 \qquad \Pr(v\in I)=\frac1{2r}\quad(v\in V(P)).
\]
\end{lemma}
\begin{proof}
Start with a stable-set distribution having inclusion probabilities at
least $1/r$. Conditional independent deletion of each included vertex
reduces every marginal to exactly $1/r$ and preserves stability.
Put $M=\lceil n/r\rceil$. Partition each resulting stable set $I$ into
parts of size at most $M$, assigning its probability mass to each part.
This gives a finite measure of total mass at most
\[
 \mathbb E\bigl(1+|I|/M\bigr)
 =1+|V(P)|/(rM)\le2.
\]
Every vertex still has mass $1/r$. Add mass on the empty set to make the
total exactly two and divide by two.
\end{proof}

We also record the standard bounded-martingale-difference estimate
\cite{azuma1967,mcdiarmid1989}, with its short exponential-moment proof.

\begin{lemma}[Bounded differences]\label{lem:concentration}
Let $Z$ be a function of $s\ge1$ independent finite random variables.
If changing one variable changes $Z$ by at most $b>0$, then
\[
 \Pr\bigl(|Z-\mathbb EZ|\ge u\bigr)
 \le2\exp\bigl(-u^2/(2sb^2)\bigr)\qquad(u>0).
\]
\end{lemma}
\begin{proof}
Expose the variables in order. The martingale of conditional
expectations has differences of conditional mean zero and absolute
value at most $b$. Convexity bounds the conditional exponential moment
of a difference at $\lambda$ by
$\cosh(\lambda b)\le\exp(\lambda^2b^2/2)$.
Multiplying these bounds and applying Markov's inequality with
$\lambda=u/(sb^2)$ gives the upper tail. Apply the same argument to $-Z$
and add the two bounds.
\end{proof}

\subsection{A small family of connectivity tests}

The next test replaces an enumeration of arbitrary vertex cuts by an
enumeration of unions of a few open neighborhoods. The family is fixed
before the random deleted set is chosen.

\begin{lemma}[Cut test]\label{lem:cut-test}
Let $P$ be a $k$-connected graph on at most $n$ vertices, where $k$ is a
sufficiently large integer. Set
\[
 s=\left\lceil10(n/k)\log(2n)\right\rceil.
\]
There is a family of at most $(n+1)^{s+1}$ matchings, each of size
$\lfloor k/8\rfloor$, with the following property. If $U\subseteq V(P)$,
every vertex of $P$ has at least $k/2$ neighbors outside $U$, and fewer
than $k/32$ endpoints of each matching belong to $U$, then $P-U$ is
$\lfloor k/32\rfloor$-connected.
\end{lemma}
\begin{proof}
For every set $R$ of at most $s$ vertices, put
$B_R=\bigcup_{v\in R}N_P(v)$. Whenever both $B_R$ and its complement
have size at least $k/3$, fix a matching of size $\lfloor k/8\rfloor$
across this cut. Such a matching exists: otherwise the endpoints of a
maximal smaller matching form a vertex cover of the crossing edges,
of size less than $k/4$. Both sides remain nonempty after their removal,
contradicting $k$-connectivity. The number of possible sets $R$, and
hence of matchings, is at most $(n+1)^{s+1}$.

The neighborhood assumption ensures that $P-U$ has more than
$\lfloor k/32\rfloor$ vertices. If it fails the asserted connectivity,
there are a set $Z\subseteq V(P)\setminus U$ with $|Z|<k/32$ and
nonempty anticomplete sides $A,A'$ partitioning $V(P)\setminus(U\cup Z)$.
Every vertex of either side has more than $k/3$ neighbors in that same
side, since it has at least $k/2$ neighbors outside $U$ and loses fewer
than $k/32$ to $Z$. In particular both sides have size at least $k/3$.

Make $s$ independent uniform draws from $A$. For each vertex $v\in A$,
the probability that none of the draws is an open neighbor of $v$ is
at most $\exp(-ks/(3n))$. The union bound over $A$ is less than one by
the choice of $s$. Thus the set $R$ of drawn vertices, after repetitions
are removed, satisfies $A\subseteq B_R$. Since $A,A'$ are anticomplete,
$A'\cap B_R=\varnothing$. The matching fixed for this cut has every edge
incident with $U\cup Z$, for otherwise an edge would join $A$ to $A'$.
But $U$ contains fewer than $k/32$ of its endpoints, and $Z$ contains
fewer than $k/32$ vertices. They cannot cover its
$\lfloor k/8\rfloor$ disjoint edges. This contradiction proves the
lemma.
\end{proof}
